
\section{Implementation}
\label{sec:implementation}

\subsection{Components Overview}

We implemented \sys on Linux 6.15, extending the kernel with a standalone file ($\sim$1 KLOC) that handles hook registration and verifier extensions. Additionally, we integrated lightweight instrumentation ($\sim$100 LOC) into NVIDIA's open GPU kernel modules to expose key memory management and scheduling decision points (\S\ref{sec:impl-host}). Our userspace loader and JIT infrastructure ($\sim$10 KLOC) builds upon bpftime's framework, while an LLVM-based backend ($\sim$1 KLOC) translates the restricted \sys eBPF instruction subset into GPU device code (PTX/SPIR-V).

\subsection{Host-side Integration}
\label{sec:impl-host}

\sys extends the NVIDIA Open GPU Kernel Modules (Turing, Ampere, Hopper, Blackwell, and newer), specifically instrumenting the UVM module at critical memory-management decision points such as the page fault handler and prefetch logic. At each hook, the driver constructs a compact context containing relevant fields (e.g., region range, fault address, block size) and invokes an attached eBPF handler. Each handler returns a small enumerated decision indicating whether to execute the default kernel logic or bypass it with custom behavior. We instrument the lifecycle events of GPU hardware queues, each corresponding to a Time-Slice Group (TSG). Specifically, we add hooks at TSG initialization (\texttt{queue_create}) and teardown (\texttt{queue_destroy}). The eBPF handlers attached to these hooks configure scheduling parameters such as TSG priority, timeslice duration, and runlist interleave frequency. Additionally, policies can cooperatively preempt active TSGs via the \texttt{gdrv_sched_preempt} kfunc, invoking the driver's existing context-switch mechanism. Our kernel module directly extends the standard Linux eBPF verifier by registering GPU-specific attach points and kfunc interfaces. The verifier extension leverages BTF annotations and kfunc type signatures to rigorously enforce type and memory safety, safely exposing critical driver-level actions such as eviction candidate selection and kernel preemption to verified eBPF handlers.


\subsection{GPU Runtime}
\label{sec:impl-gpu}

\paragraph{Code Generation and Binary Trampolines.}

Once a \sys eBPF program passes verification and is loaded into the kernel, our loader intercepts its bytecode and BTF metadata, then forwards them to our userspace backend for GPU code generation (PTX/SPIR-V). Verified code executes either directly in the kernel (using the standard eBPF JIT) or on the GPU through our specialized backend and runtime. To support runtime instrumentation without recompilation overhead, \sys dynamically injects lightweight binary trampolines at strategic kernel points (entry points, selected memory instructions, and execution-phase boundaries). We also expose a minimal set of runtime helper functions, pre-compiled to GPU device code with verifier-approved signatures. During translation, static analysis ensures minimal context structures are constructed, only materializing fields referenced by handlers, and inserts prologue/epilogue logic to manage SIMT-aware register state preservation. By directly translating verified eBPF bytecode to GPU-executable code, \sys seamlessly extends Linux's kernel-level safety guarantees (e.g. memory safety, bounded loops, and type safety) into GPU kernels.

\paragraph{Scheduling and Memory Operations.}

For scheduling policies, \sys employs a thread-block scheduler abstraction that models kernel execution as dynamically scheduled units of work (e.g., tiles or logical blocks), executed by persistent worker blocks. On Blackwell GPUs, we leverage hardware-supported cluster launch control APIs, which treat block grids as dynamically manageable tasks, enabling efficient hardware-supported work-stealing. Each GPU worker executes its assigned unit, invokes device-side \sys policy handlers (running in dedicated SIMT lanes) to select or approve subsequent tasks, and interacts with the cluster control APIs to claim or relinquish tasks. Execution-phase annotations and instrumentation points are implemented via low-overhead \sys trampolines. Critically, all GPU-resident scheduling logic strictly adheres to bounded-execution constraints enforced by our verifier.

For memory management, \sys attaches handlers to select memory instructions, device-side tracepoints, or kernel function boundaries, allowing policies to issue lightweight hints back to the host driver through helper calls. In contrast to prior GPU instrumentation efforts (e.g., eGPU) that rely on expensive atomic synchronization and fences, \sys helper functions and instrumentation logic avoid explicit synchronization primitives, significantly minimizing overhead.

\paragraph{Cross-Domain State and Synchronization.}

\sys relies on standard eBPF maps as its unified state abstraction across host-driver and GPU-device boundaries. To ensure efficient, zero-copy synchronization, we allocate map structures within GPU-accessible unified memory regions, thereby eliminating explicit data marshalling overhead between CPU and GPU contexts. These shared maps maintain dynamic policy state such as per-region access frequency (hotness), per-tenant scheduling priorities, block-level execution progress, and instrumentation counters. Unlike previous eGPU implementations, which restricted maps exclusively to CPU memory, causing prohibitive synchronization overhead, \sys transparently supports placing frequently accessed policy state directly within GPU registers, GPU local memory, or CPU-side memory, enabling optimal locality and minimal runtime overhead.

\subsection{SIMT-aware Verification Implementation}
\label{sec:impl-verification}

The host-side verifier uses Linux kernel's existing eBPF verifier without modification, ensuring standard memory safety, bounded loops, and resource usage. We implement our GPU-specific extensions by adding an extra verification pass at program load time within the kernel module. This pass statically identifies warp-uniform fields from context structures annotated via BTF metadata. The verifier then performs conservative dataflow analysis: each register and stack slot is tagged as uniform or varying, and these tags propagate through arithmetic and logical operations. Combining two uniform values yields uniform; any operation involving a varying operand yields varying; helper return values use uniformity annotations from their BTF declarations. Control flow analysis checks that branches and loops depend solely on these warp-uniform fields, explicitly rejecting programs violating this constraint. Map update operations are additionally checked to ensure keys are warp-uniform.

We enforce resource bounds per hook type, setting conservative instruction and memory-operation budgets based on the expected trigger frequency and latency budget of each hook. Our verifier explicitly rejects GPU-wide synchronization constructs, global barriers, and non-uniform atomic operations that could cause hardware stalls or memory consistency violations. No runtime sandboxing mechanisms are used; policies that pass verification run directly on GPU hardware, while rejected programs are never executed.

\subsection{Warp-level Execution Implementation}

The GPU runtime implements warp-level execution aggregation via lightweight device-side trampolines inserted at kernel entry, synchronization barriers, and critical memory instructions. The two-phase execution model is realized as follows: in the lane-local phase, each lane computes its per-thread contributions (e.g., effective addresses, partial counters) into registers; these are then aggregated using \texttt{\_\_ballot\_sync} for masks and \texttt{\_\_shfl\_sync} for value reduction. In the warp-policy phase, lane~0 (or the first active lane via \texttt{\_\_ffs(\_\_activemask())}) executes the translated eBPF handler with the aggregated inputs. Trampolines minimally save warp execution state, aggregate lane-varying inputs using GPU shuffle instructions, execute the scalar policy handler once per warp, and then broadcast results back to all threads using shuffle operations.

We implement aggressive function and map-helper call inlining at JIT compilation, significantly reducing call overhead and branching. GPU map operations are specialized by identifying common access patterns and embedding them directly into the handler code, reducing register pressure and avoiding global memory round trips. An optional sampling mechanism is also provided to allow dynamic control of policy overhead.

\subsection{Hierarchical Map Implementation}

Our runtime internally employs a cost-based heuristic placement model for cross-layer maps. Low-frequency global state is stored in CPU host DRAM with periodic GPU-side snapshots cached locally. Frequently accessed GPU-intensive state is primarily stored in GPU global memory, periodically aggregated and asynchronously flushed back to host memory for policy consumption. Extremely hot per-warp or per-SM data resides directly in GPU shared memory or SM-local storage, ensuring extremely low latency access.

Device-side handlers use standard eBPF helper APIs to access maps, fully unaware of underlying memory management complexity. Per-key atomicity is implemented using native GPU atomic operations for single-key updates. The snapshot-based consistency model is realized by defining merge points at GPU kernel completion boundaries: when a kernel finishes, the runtime triggers an asynchronous flush of GPU-local map shards to canonical instances in global memory or host DRAM. Host-side handlers and user-space control-plane logic observe these coherent snapshots rather than in-flight GPU-side updates.

As a concrete example, GPU-side memory-access instrumentation handlers record per-region statistics directly into GPU-local map shards. Periodically, the runtime aggregates these shards into global GPU memory, subsequently making them available to host-side handlers. Host-side eviction handlers then read coherent aggregate statistics directly from global maps, enabling accurate memory management decisions without explicit data movement or synchronization logic in policy handlers.
